\section{点涡平衡点处的稳定性分析}
\subsection{控制方程}
\subsubsection{点涡速度控制方程（原始形式）}
\begin{align}
    \frac{\mathrm{d}z_{0}}{\mathrm{d}t}=\frac{\Gamma}{2\pi i}\left( \frac{c-1}{2z_{0}} - \frac{cz_{0}^{c-1}}{z_{0}^{c}-\bar{z}_{0}^{c}} \right) - Scz_{0}^{c-1}
\end{align}

\subsubsection{无量纲化过程}

引入特征尺度：
\[
z_0 = L \tilde{z}_0, \quad t = T \tilde{t}
\]

代入后得：
\begin{align*}
    \frac{\mathrm{d}\tilde{z}_{0}}{\mathrm{d}\tilde{t}} &= \frac{T}{L} \cdot \frac{\Gamma}{2\pi i} \left( \frac{1}{L} \cdot \frac{c-1}{2\tilde{z}_0} - \frac{1}{L} \cdot \frac{c\tilde{z}_0^{c-1}}{\tilde{z}_0^c - \bar{\tilde{z}}_0^c} \right) - \frac{T}{L} SL^{c-1} c\tilde{z}_0^{c-1}
\end{align*}

令系数归一化：
\begin{align*}
    \frac{T\Gamma}{2\pi L^{2}}=1, \quad TSL^{c-2}=1
\end{align*}

解得：
\begin{align*}
    L = \left( \frac{\Gamma}{2\pi S} \right)^{1/c}, \quad
    T = S^{-1} \left( \frac{\Gamma}{2\pi S} \right)^{\frac{2}{c} - 1}
\end{align*}

\subsubsection{无量纲点涡控制方程}

\begin{align}
    \boxed{
    \frac{\mathrm{d}\tilde{z}_{0}}{\mathrm{d}\tilde{t}} = \frac{1}{i} \left( \frac{c-1}{2\tilde{z}_0} - \frac{c\tilde{z}_0^{c-1}}{\tilde{z}_0^c - \bar{\tilde{z}}_0^c} \right) - c \tilde{z}_0^{c-1}
    }
\end{align}

\subsubsection{无量纲流场复速度表达式}

\begin{align}
    \frac{\mathrm{d}\tilde{W}}{\mathrm{d}\tilde{z}} = \frac{c}{i} \cdot \tilde{z}^{c-1} \left( \frac{1}{\tilde{z}^{c} - \tilde{z}_0^{c}} - \frac{1}{\tilde{z}^{c} - \bar{\tilde{z}}_0^{c}} \right) - c \tilde{z}^{c-1}
\end{align}

流体质点在复速度诱导下的路径为：
\begin{align}
    \frac{d\tilde{z}}{d\tilde{t}} = \left. \frac{d\tilde{W}}{d\tilde{z}} \right|_{\tilde{z}}
\end{align}

\subsection{点涡的运动分析}
设点涡的实际位置坐标$z_{0}$，其在背景流中保持静止时位置坐标$z_{s}$，则有：
\begin{align}
    z_{0}=z_{s}+\Delta z_{0}
\end{align}
其中 $\Delta z_0$ 表示点涡相对于平衡位置的微小扰动。

前文构造的复势结构是通过令 $z_0 = z_s$ 并满足 $\left.\frac{dW}{dz}\right|_{z_s} = 0$ 来确定平衡状态。
在该基础上，为研究点涡在扰动下的动态行为，可引入扰动项 $\Delta z_0$ 并开展稳定性分析。

点涡位置的随时间演化由背景复势诱导速度所决定，其动力学方程为：
\begin{align}
    \frac{\mathrm{d}z_{0}}{\mathrm{d}t}=\left.\frac{\mathrm{d}W}{\mathrm{d}z}\right|_{z=z_{0}}=\frac{\Gamma}{2\pi i}\left(\frac{c-1}{2z_{0}}-\frac{cz_{0}^{c-1}}{z_{0}^{c}-\bar{z}_{0}^{c}}\right)-Scz_{0}^{c-1}
\end{align}
上式中的复速度表达式已经在奇异性分析的基础上，剔除了点涡自身的对数奇性项，所保留的$\frac{c-1}{2z_{0}}$正是该自感应速度在展开后的有限残差。
因此，该表达式可用于描述点涡在镜像点涡与背景流共同作用下的运动行为，进而用于平衡态构造与稳定性分析

点涡作为背景复势诱导速度场中的运动粒子，其速度由其余流场诱导决定。若要使其保持静止，必须满足：
\begin{align}
    \left.\frac{\mathrm{d}z_{0}}{\mathrm{d}t}\right|_{z_{0}=z_{s}}=\left.\frac{\mathrm{d}W}{\mathrm{d}z}\right|_{z=z_{s}}=0
\end{align}
此时的 $z_s$ 即为通过静止条件解得的点涡平衡位置，代表一个由势流系统自洽确定的特殊结构点。

\subsubsection{控制方程的线性化}
设点涡位置坐标及共轭位置坐标在稳定点$z_{s}$附近展开为
\begin{align*}
    z_{0}&=z_{s}+\Delta z_{0}\\
    \bar{z}_{0}&=\bar{z}_{s}+\Delta \bar{z}_{0}
\end{align*}

记无量纲点涡运动控制方程为
\begin{align*}
    \frac{\mathrm{d}{z}_{0}}{\mathrm{d}{t}} &=f(z_{0},\bar{z}_{0})\\
    &=\frac{1}{i} \left( \frac{c-1}{2{z}_0} - \frac{c{z}_0^{c-1}}{{z}_0^c - \bar{{z}}_0^c} \right) - c {z}_0^{c-1}\\
    &=\frac{c-1}{2i}\frac{1}{z_{0}}-\frac{c}{i}\frac{z_{0}^{c-1}}{z_{0}^{c}-\bar{z}_{0}^{c}}-cz_{0}^{c-1}
\end{align*}
将点涡运动的扰动形式代入点涡控制方程，并在$z_{s}$处一阶Taylor展开
\begin{align*}
    \frac{\mathrm{d}}{\mathrm{d}t}(z_{s}+\Delta z_{0})=f(z_{s})+f'(z_{s})\Delta z_{0}+\cdots
\end{align*}
由于$z_{s}$为静止点，故$f(z_{s})=0$，保留方程的线性项：
\begin{align*}
    \frac{\mathrm{d}}{\mathrm{d}t}(\Delta z)&\approx \left.\frac{\partial f}{\partial z_{0}}\right|_{z_{s}}\Delta z_{0}+\left.\frac{\partial f}{\partial \bar{z}_{0}}\right|_{\bar{z}_{s}}\Delta \bar{z}_{0}\\
    &=a\Delta z_{0}+b\Delta \bar{z}_{0}
\end{align*}
其中$a,b$为复常数，可以记为$a=a_{r}+ia_{i},b=b_{r}+ib_{i}$。

求偏导
\begin{align*}
    \frac{\partial f}{\partial z_{0}}&=\frac{1-c}{2i}\frac{1}{z_{0}^{2}}-\frac{c}{i}\frac{z_{0}^{c-2}[(1-c)\bar{z}_{0}^{c}-z_{0}^{c}]}{(z_{0}^{c}-\bar{z}_{0}^{c})^{2}}-c(c-1)z_{0}^{c-2}\\
    \frac{\partial f}{\partial \bar{z}_{0}}&=-\frac{c}{i}\frac{c(z_{0}\bar{z}_{0})^{c-1}}{(z_{0}^{c}-\bar{z}_{0}^{c})^{2}}
\end{align*}
则
\begin{align*}
    a&=\frac{1-c}{2i}\frac{1}{z_{s}^{2}}-\frac{c}{i}\frac{z_{s}^{c-2}[(1-c)\bar{z}_{s}^{c}-z_{s}^{c}]}{(z_{s}^{c}-\bar{z}_{s}^{c})^{2}}-c(c-1)z_{s}^{c-2}\\
    b&=-\frac{c}{i}\frac{c(z_{s}\bar{z}_{s})^{c-1}}{(z_{s}^{c}-\bar{z}_{s}^{c})^{2}}
\end{align*}

设$\Delta z=x+iy$，将复变量微分系统转化为二维实变量微分系统
\begin{align*}
    \frac{\mathrm{d}}{\mathrm{d}t}(x+iy)&=(a_{r}+ia_{i})(x+iy)+(b_{r}+ib_{i})(x-iy)\\
&=[(a_{r}+b_{r})x+(b_{i}-a_{i})y]+i[(a_{i}+b_{i})x+(a_{r}-b_{r})y]\\
\end{align*}
复微分方程等价于
\begin{align*}
    \frac{\mathrm{d}x}{\mathrm{d}t}&=(a_{r}+b_{r})x+(b_{i}-a_{i})y\\
    \frac{\mathrm{d}y}{\mathrm{d}t}&=(a_{i}+b_{i})x+(a_{r}-b_{r})y
\end{align*}
写成矩阵形式
\begin{align*}
    \frac{\mathrm{d}}{\mathrm{d}t}
\begin{bmatrix}
    x \\
    y 
\end{bmatrix}
&=
\begin{bmatrix}
    a_{r}+b_{r} & b_{i}-a_{i} \\
    a_{i}+b_{i} & a_{r}-b_{r}
\end{bmatrix}
\begin{bmatrix}
    x \\
    y 
\end{bmatrix}\\
&=
\begin{bmatrix}
    \Re(a+b) & -\Im(a-b) \\
    \Im(a+b) &  \Re(a-b)
\end{bmatrix}
\begin{bmatrix}
    x \\
    y 
\end{bmatrix}
\end{align*}
该系统的Jacobi矩阵结构是由于复变量的共轭变量不可省略，因此无法像解析单变量系统一样通过 $f'(z)$ 导数简化

\subsubsection{直角区域的点涡运动分析}
取$c=2$，点涡位置在直角区域随时间的演化动力学方程为
\begin{align*}
    \frac{\mathrm{d}z_{0}}{\mathrm{d}t}=\frac{i\Gamma}{2\pi}\left(\frac{2z_{0}}{z_{0}^{2}-\bar{z}_{0}^{2}}-\frac{1}{2z_{0}} \right) - 2Sz_{0}
\end{align*}
取对应的特征尺度:
\begin{align*}
    L = \frac{1}{2S}, \quad
    T = \frac{\pi}{2S^{2}\Gamma}
\end{align*}
直角区域中无量纲点涡控制方程
\begin{align}
    \boxed{
    \frac{\mathrm{d}\tilde{z}_{0}}{\mathrm{d}\tilde{t}} = i\left( \frac{2\tilde{z}_0}{\tilde{z}_0^2 - \bar{\tilde{z}}_0^2} -\frac{1}{2\tilde{z}_0} \right) - \tilde{z}_0
    }
\end{align} 
无量纲流场复速度表达式
\begin{align*}
    \frac{\mathrm{d}\tilde{W}}{\mathrm{d}\tilde{z}} = \frac{2S\Gamma}{\pi i} \cdot \tilde{z} \left( \frac{1}{\tilde{z}^{2} - \tilde{z}_0^{2}} - \frac{1}{\tilde{z}^{2} - \bar{\tilde{z}}_0^{2}} \right) - \tilde{z}
\end{align*}
若要使点涡在流场中保持静止，必须满足：
\begin{align}
    \left.\frac{\mathrm{d}z_{0}}{\mathrm{d}t}\right|_{z_{0}=z_{s}}=\left.\frac{\mathrm{d}W}{\mathrm{d}z}\right|_{z=z_{s}}=0
\end{align}
此时的 $z_s$ 即为通过静止条件解得的点涡平衡位置，代表一个由势流系统自洽确定的特殊结构点。

将点涡动力学方程线性化，记
\begin{align}
    \frac{\mathrm{d}{z}_{0}}{\mathrm{d}{t}} &=f(z_{0},\bar{z}_{0}) \\
    &=i\left( \frac{2{z}_0}{{z}_0^2 - \bar{{z}}_0^2} -\frac{1}{2{z}_0} \right) - {z}_0
\end{align}
由于$z_{s}$为静止点，故$f(z_{s})=0$，保留方程的线性项：
\begin{align*}
    \frac{\mathrm{d}}{\mathrm{d}t}(\Delta z)&\approx \left.\frac{\partial f}{\partial z_{0}}\right|_{z_{s}}\Delta z_{0}+\left.\frac{\partial f}{\partial \bar{z}_{0}}\right|_{\bar{z}_{s}}\Delta \bar{z}_{0}\\
    &=a\Delta z_{0}+b\Delta \bar{z}_{0}
\end{align*}
其中$a,b$为复常数，可以记为$a=a_{r}+ia_{i},b=b_{r}+ib_{i}$。

求偏导数
\begin{align*}
    \frac{\partial f}{\partial z_{0}}&=
    i\left[\frac{1}{2z_{0}^{2}}-\frac{1}{(z_{0}+\bar{z}_{0})^{2}}-\frac{1}{(z_{0}-\bar{z}_{0})^{2}}\right]-1\\
    \frac{\partial f}{\partial \bar{z}_{0}}&=
    i\left[\frac{1}{(z_{0}-\bar{z}_{0})^{2}}-\frac{1}{(z_{0}+\bar{z}_{0})^{2}}\right]
\end{align*}
解得$a,b$
\begin{align*}
    a&=
    i\left[\frac{1}{2z_{s}^{2}}-\frac{1}{(z_{s}+\bar{z}_{s})^{2}}-\frac{1}{(z_{s}-\bar{z}_{s})^{2}}\right]-1\\
    b&=i\left[\frac{1}{(z_{s}-\bar{z}_{s})^{2}}-\frac{1}{(z_{s}+\bar{z}_{s})^{2}}\right]
\end{align*}

静止点 $z_{s}$ 的表达式为
\begin{align}
    \boxed{
    z_{s} = \left( \frac{iS\Gamma}{2\pi } \right)^{1/2}
    = \frac{1}{\sqrt{2}} \sqrt{\frac{S\Gamma}{\pi}} (1 + i)
    }
\end{align}

推导如下：我们将 \( i = e^{i\frac{\pi}{2}} \)，因此有
\begin{align*}
    z_s 
    &= \left( \frac{iS\Gamma}{2\pi} \right)^{1/2}
    = \left( \frac{S\Gamma}{2\pi} \cdot e^{i\frac{\pi}{2}} \right)^{1/2} \\
    &= \sqrt{\frac{S\Gamma}{2\pi}} \cdot e^{i\frac{\pi}{4}}
    = \sqrt{\frac{S\Gamma}{2\pi}} \cdot \left( \cos\frac{\pi}{4} + i\sin\frac{\pi}{4} \right) \\
    &= \sqrt{\frac{S\Gamma}{2\pi}} \cdot \left( \frac{1}{\sqrt{2}} + i\cdot \frac{1}{\sqrt{2}} \right)
    = \frac{1}{\sqrt{2}} \sqrt{\frac{S\Gamma}{\pi}} (1 + i)
\end{align*}

因此，静止点 $z_s$ 具有实部和虚部相等的复数结构，位于第一象限，代表点涡在二维空间中可能实现稳定平衡的一个特征性位置。

\textbf{注：}此处所使用的 $z_s$ 已经过无量纲化处理，实际上反映的是以特征长度 $L = 1/(2S)$ 为尺度归一化后的坐标结构。

这正是点涡在合适的背景线性剪切流中产生的一个稳定平衡点，其存在的条件是：

\begin{itemize}
    \item 点涡诱导的速度与背景流在该点的剪切速度相互平衡；
    \item 背景剪切强度 $S$ 与点涡强度 $\Gamma$ 恰好匹配；
    \item 而这个“匹配关系”正由 $z_s$ 的坐标结构确定。
\end{itemize}

我们将 $z_s = \left( \frac{iS\Gamma}{2\pi} \right)^{1/2}$ 写为指数形式：
\begin{align*}
    z_s &= \left( \frac{S\Gamma}{2\pi} \right)^{1/2} \cdot e^{i\pi/4} = \frac{1}{\sqrt{2}} \sqrt{ \frac{S\Gamma}{\pi} } (1+i)
\end{align*}
这表明：当点涡位于第一象限、满足 $\Re(z_s) = \Im(z_s)$ 的结构时，诱导速度场与背景剪切场在该点方向一致、大小相等，从而形成自洽平衡。

从物理上来看，该平衡点的形成是在边界条件约束下，通过点涡自诱导速度与背景剪切速度的平衡而实现的。边界的存在通过复势函数中的镜像结构和映射形式已经隐含地施加在动力学系统中。因此，该平衡点是边界控制下的“自诱导-外剪切平衡”的典型产物，体现了点涡系统与壁面之间复杂而精确的相互作用。

设$\Delta z=x+iy$，将复变量微分系统转化为二维实变量微分系统并写为矩阵形式

\begin{align*}
    \frac{\mathrm{d}}{\mathrm{d}t}
\begin{bmatrix}
    x \\
    y 
\end{bmatrix}
&=
\begin{bmatrix}
    \Re(a+b) & -\Im(a-b) \\
    \Im(a+b) &  \Re(a-b)
\end{bmatrix}
\begin{bmatrix}
    x \\
    y 
\end{bmatrix}
\end{align*}
由于点涡所处背景流场包含复共轭结构，扰动系统在$z_{0},\bar{z}_{0}$方向均具有非零耦合，因此线性系统表现为非Hermitian耦合形式，Jacobian矩阵依赖于$a,b$ 的复实部组合。

\subsubsection{线性扰动系统的通解}

记线性系统为
\begin{align*}
    \frac{\mathrm{d}\bm{X}}{\mathrm{d}t} = \bm{JX}
\end{align*}
其通解形式为
\begin{align*}
    \bm{X}(t) = e^{\bm{J}t} \bm{X}_{0}
\end{align*}
其中 $\bm{X}_0 = \begin{bmatrix} x_0 \\ y_0 \end{bmatrix}$ 为初始扰动，$\bm{J}$ 为实系数矩阵。我们通过特征值–特征向量方法构造 $e^{\bm{J}t}$。

根据前述分析，$\bm{J}$ 的结构由平衡点 $z_s$ 处的偏导 $a,b$ 决定。为分析简便，直接指定
\[
z_s = \frac{1}{2} + i\frac{1}{2}
\]
代入得
\[
a = 0, \quad b = -2i
\]
从而得到
\[
\bm{J} =
\begin{bmatrix}
    \Re(a + b) & -\Im(a - b) \\
    \Im(a + b) & \Re(a - b)
\end{bmatrix}
=
\begin{bmatrix}
    0 & -2 \\
    -2 & 0
\end{bmatrix}
\]

\paragraph{特征值与特征向量}
矩阵的特征方程为
\begin{align*}
    \det(\bm{J} - \lambda \bm{I}) &= 
    \begin{vmatrix}
        -\lambda & -2 \\
        -2 & -\lambda
    \end{vmatrix}
    = \lambda^2 - 4 = 0
\end{align*}
解得特征值
\[
\lambda = \pm 2
\]

对应特征向量为：
\begin{align*}
\lambda = 2 \quad &\Rightarrow\quad \bm{v}_1 = \begin{bmatrix} 1 \\ -1 \end{bmatrix} \\
\lambda = -2 \quad &\Rightarrow\quad \bm{v}_2 = \begin{bmatrix} 1 \\ 1 \end{bmatrix}
\end{align*}

\paragraph{通解表达}
\begin{align*}
\begin{bmatrix}
    x(t) \\
    y(t)
\end{bmatrix}
=
C_1 e^{2t}
\begin{bmatrix}
    1 \\
    -1
\end{bmatrix}
+
C_2 e^{-2t}
\begin{bmatrix}
    1 \\
    1
\end{bmatrix}
\end{align*}

\paragraph{转为复变量形式}
将实变量表示转化为复变量 $\Delta z_0(t) = x(t) + i y(t)$，得：
\begin{align*}
    \Delta z_0(t) &= C_1 e^{2t} (1 - i) + C_2 e^{-2t} (1 + i)
\end{align*}

记
\begin{align*}
    \Delta z_0(t) = \epsilon_1 e^{2t} + \epsilon_2 e^{-2t}, \quad
    \epsilon_1 = C_1(1 - i), \quad
    \epsilon_2 = C_2(1 + i)
\end{align*}
\subsubsection{引入物理约束下的扰动解结构}

通解可写为：
\begin{align*}
    \Delta z_0(t) = \epsilon_1 e^{2t} + \epsilon_2 e^{-2t}, \quad
    \epsilon_1 = C_1(1 - i), \quad \epsilon_2 = C_2(1 + i)
\end{align*}

根据物理背景，引入如下约束：

\begin{itemize}
    \item \textbf{稳定性要求}：扰动不能发散，必须令 $\epsilon_1 = 0$，即 $C_1 = 0$；
    \item \textbf{扰动旋转性}：希望扰动以某固定角速度绕稳定点旋转；
    \item \textbf{扰动收敛性}：扰动应随时间衰减回稳定点，即模长随时间单调减小；
\end{itemize}

因此扰动形式简化为：
\begin{align*}
    \Delta z_0(t) = \epsilon_2 e^{-2t}
\end{align*}

其中 $\epsilon_2 = C_2(1 + i)$ 是复系数，控制扰动初始相位与旋转方向。

该扰动为螺旋状收敛轨迹，其模长 $|\Delta z_0(t)| = |\epsilon_2| e^{-2t}$，随时间指数衰减，体现了线性稳定性。
当 $\epsilon_2 = \epsilon e^{-2it}$ 时，为纯旋转结构，需满足更严格条件 $b=0$，此为特例。

因此，本扰动模型从通解出发，在满足物理稳定性与旋转需求的前提下，构建出具有明确物理含义的解。


我们在线性扰动框架下重构了点涡在直角势流背景中演化的动力学行为。通过在平衡点处对演化方程进行复变量线性化，我们导出了一个包含两维实变量的线性系统通解，并明确指出论文中给出的旋转扰动解 $\Delta z(t) = \epsilon e^{-2it}$ 实际为该通解在特定初值条件（扰动模长不变）下的一个特解。物理上，该扰动模式对应点涡围绕平衡点以固定半径匀速旋转的稳定状态，其旋转频率为 2，周期为 $\pi$，与原文实验或理论结论一致。这一分析不仅验证了论文中的推导过程，也展示了通用线性扰动分析在势流-涡旋耦合问题中的适用性与解释力。

\subsection{从物理参数出发的扰动稳定性分析}

本节我们从背景剪切强度 \( S \) 和点涡强度 \( \Gamma \) 出发，构造完整的扰动演化分析流程。点涡的无量纲动力学方程为：
\begin{align}
    \frac{\mathrm{d}z_0}{\mathrm{d}t} = i\left( \frac{2z_0}{z_0^2 - \bar{z}_0^2} - \frac{1}{2z_0} \right) - z_0
\end{align}

\paragraph{静止点的确定}

设点涡处于平衡状态，即：
\begin{align}
    \left. \frac{\mathrm{d}z_0}{\mathrm{d}t} \right|_{z_0 = z_s} = 0
\end{align}
代入上式，得到静止点满足的代数方程：
\begin{align}
    i\left( \frac{2z_s}{z_s^2 - \bar{z}_s^2} - \frac{1}{2z_s} \right) - z_s = 0
\end{align}
该方程通常为复代数方程，可以通过数值方式求解所有满足条件的 \( z_s \)。

\paragraph{线性扰动控制系统}

记扰动形式 \( \Delta z = x + i y \)，则系统可线性化为：
\begin{align}
    \frac{\mathrm{d}}{\mathrm{d}t}
    \begin{bmatrix}
        x \\ y
    \end{bmatrix}
    =
    \bm{J}
    \begin{bmatrix}
        x \\ y
    \end{bmatrix}, \quad
    \bm{J} =
    \begin{bmatrix}
        \Re(a + b) & -\Im(a - b) \\
        \Im(a + b) & \Re(a - b)
    \end{bmatrix}
\end{align}
其中：
\begin{align*}
    a &= i\left[\frac{1}{2z_s^2} - \frac{1}{(z_s + \bar{z}_s)^2} - \frac{1}{(z_s - \bar{z}_s)^2}\right] - 1 \\
    b &= i\left[\frac{1}{(z_s - \bar{z}_s)^2} - \frac{1}{(z_s + \bar{z}_s)^2} \right]
\end{align*}

\paragraph{扰动通解及稳定性判据}

求解线性系统特征值：
\[
\det(\bm{J} - \lambda \bm{I}) = 0
\]
记特征值为 \( \lambda_{1,2} \)，对应特征向量为 \( \bm{v}_{1,2} \)，则扰动演化的通解为：
\begin{align}
    \begin{bmatrix}
        x(t) \\ y(t)
    \end{bmatrix}
    = C_1 e^{\lambda_1 t} \bm{v}_1 + C_2 e^{\lambda_2 t} \bm{v}_2
\end{align}
转回复变量：
\begin{align}
    \Delta z_0(t) = x(t) + i y(t)
\end{align}

\paragraph{物理约束与解的选择}

在物理问题中，我们通常要求扰动满足以下条件：
\begin{itemize}
    \item 扰动不能发散：排除 \( \Re(\lambda) > 0 \) 的模式；
    \item 保持幅度有限：选择指数衰减或纯旋转形式；
    \item 指定初始扰动方向：由边界条件或初始流场决定。
\end{itemize}


设排除发散项，即令 \( C_1 = 0 \)，则解为：
\begin{align}
    \Delta z_0(t) = C_2 e^{\lambda_2 t} (v_{2x} + i v_{2y})
\end{align}
其中 \( \lambda_2 \) 是具有负实部或纯虚数的特征值。

\paragraph{关于多个平衡点的说明}

由于复代数方程可能具有多个根，系统存在多个静止点 \( \{ z_s^{(1)}, z_s^{(2)}, \ldots \} \)，每一个都对应一个线性扰动系统 \( \bm{J}^{(k)} \)。

处理方法：
\begin{itemize}
    \item 分别分析每一个静止点的稳定性；
    \item 根据物理设置（对称性、边界形状）选择实际存在的点；
    \item 用全局数值模拟验证各点的吸引/排斥性质。
\end{itemize}



